% =========================================================
% Proof file: prf-minus-one-times-vector.tex
% =========================================================

\newpage
\phantomsection
\label{prf:minus-one-times-vector}

\begin{remark*}[Return]
\hyperref[thm:minus-one-times-vector]{Return to Theorem}
\end{remark*}

\begin{theorem*}[($-1)v=-v$]
For every $v \in V$,
\[
(-1)v=-v.
\]
\end{theorem*}

\begin{proof}
\textbf{Professional Standard Proof.}~\\
Let $v \in V$. Then
\[
v+(-1)v=1v+(-1)v.
\]
By distributivity,
\[
1v+(-1)v=(1+(-1))v=0v.
\]
By \hyperref[thm:zero-scalar-times-vector]{$0v=0$}, it follows that
\[
v+(-1)v=0.
\]
Thus $(-1)v$ is an additive inverse of $v$. By uniqueness of additive inverses,
\[
(-1)v=-v.
\]
\end{proof}

\begin{proof}
\textbf{Detailed Learning Proof.}~\\

\textbf{Step 1.} Fix an arbitrary vector.\\
Let $v \in V$. The goal is to show that scalar multiplication by $-1$ produces the additive inverse of $v$.

\textbf{Step 2.} Rewrite $v$ as $1v$.\\
By the scalar identity axiom,
\[
1v=v.
\]
Therefore
\[
v+(-1)v=1v+(-1)v.
\]

\textbf{Step 3.} Use distributivity over scalar addition.\\
Applying distributivity gives
\[
1v+(-1)v=(1+(-1))v.
\]
Because $1+(-1)=0$ in $\mathbf{F}$, this becomes
\[
1v+(-1)v=0v.
\]

\textbf{Step 4.} Use the theorem $0v=0$.\\
By \hyperref[thm:zero-scalar-times-vector]{$0v=0$},
\[
0v=0.
\]
Hence
\[
v+(-1)v=0.
\]

\textbf{Step 5.} Interpret the resulting equation.\\
The equation
\[
v+(-1)v=0
\]
shows that $(-1)v$ is an additive inverse of $v$.

\textbf{Step 6.} Use uniqueness of additive inverses.\\
The vector $-v$ is defined to be the unique additive inverse of $v$. Since $(-1)v$ is also an additive inverse of $v$, uniqueness implies
\[
(-1)v=-v.
\]

\textbf{Step 7.} State the conclusion.\\
Thus scalar multiplication by $-1$ yields the additive inverse of a vector.
\end{proof}

\begin{remark*}[Proof structure]
The proof is direct. It shows first that $(-1)v$ satisfies the defining property of an additive inverse, and then invokes uniqueness of additive inverses.
\end{remark*}

\begin{remark*}[Dependencies]
\hfill
\begin{itemize}
  \item \hyperref[def:vector-space]{Definition of Vector Space}
  \item \hyperref[thm:zero-scalar-times-vector]{$0v=0$}
  \item \hyperref[thm:unique-additive-inverse]{Unique Additive Inverse}
\end{itemize}
\end{remark*}